Streamers are elongated structures within protostellar envelopes that connect large radii to the forming disks, as revealed by high-resolution (sub-)millimeter interferometric observations. These structures have been identified as potential channels for transporting a non-negligible fraction of envelope mass onto protoplanetary disks, which may influence disk growth and dynamics. Therefore, quantifying the origin and kinematics of streamers is essential for understanding the mechanisms by which forming protoplanetary disks acquire mass from their surrounding envelopes.

In this thesis, we develop a dynamical model in which streamers arise from localized density enhancements within the collapsing envelope, motivated by the growth of gravitational instabilities. Observationally, such overdensities are traced indirectly as enhancements in line emission in moment maps and position--position--velocity (PPV) data cubes. We test the model against four systems (Per-emb-2 and Per-emb-50 observed with NOEMA, and S~CrA and HL~Tau observed with ALMA) by fitting the projected morphology and line-of-sight velocity information. Specifically, we constrain the model parameters ($\theta$, $\phi_0$, $i$, $t_{\rm s}$, $\omega$) using point-method representation derived from observation data and global-method that compares the model with the full position-position-velocity (PPV) data cube.

Overall, the proposed model provides a good match to the observed streamer morphology and kinematics across all targets, demonstrating that the essential features of the data can be reproduced consistently. The inferred model parameters allow us to derive key envelope properties, including the three-dimensional streamer orientation, the underlying envelope density profile, and the associated mass accretion rates. Our results indicate that, although streamers represent coherent accretion channels linking the envelope to disk-forming scales, their contribution to the total mass accretion budget onto the protoplanetary disk is generally subdominant. This suggests that streamers play a secondary role in disk growth, while still serving as important tracers of non-axisymmetric infall in collapsing protostellar envelopes.

\vspace{1em}
\noindent\textbf{Keywords:} streamers; protostellar core; star formation; embedded protoplanetary disk; high-resolution radio interferometry